\subsection{The obstruction survives modifications of the base}

\begin{proposition}[B.4.27; Proved; no log-torsor repair by point blowups]
\label{push:B427}
Let $B=\Delta_Q\times\Delta_R$, with
$q=Q^{12m}$ and $r=R^n$, where $m,n$ are positive
integers. Let $\beta:S\to B$ be any finite sequence
of point blowups whose centres lie over the origin.
Give $S$ the divisorial log structure of the reduced
inverse image of $QR=0$, and let $\mathcal A_S$ be
the pullback of the actual Jacobian log torus of
B.4.24. Then
\[
 H^1(S,\mathcal A_S)=0
 \tag{B.4.27.1}
\]
for the usual analytic topology. In particular the
actual associated axial log-torus torsor on the
punctured bidisc has no locally trivial
$\mathcal A_S$-torsor extension after these
modifications.

This result concerns the specified associated
log-torus sheaf. It does not assert that flattening
or resolving an arbitrary analytic origin cap
produces a torsor under that sheaf.
\end{proposition}

\begin{proof}
Write $E_1,\ldots,E_N$ for the exceptional curves,
and $D_Q,D_R$ for the strict transforms of the
two coordinate axes. The reduced total boundary
is normal crossing. Every $E_i$ is a rational
curve, and the exceptional divisor is a connected
tree if $N>0$. Its two orders
\[
 a_i=\operatorname{ord}_{E_i}(Q),\qquad
 b_i=\operatorname{ord}_{E_i}(R)
 \tag{B.4.27.2}
\]
are positive: the image of $E_i$ is the origin.
Set
\[
 D_1=D_Q+\sum_i E_i,\qquad
 D_2=D_Q+D_R+\sum_i E_i.
\]
For a reduced normal-crossing divisor $D$, denote
by $\mathcal M_D^{\mathrm{gp}}$ the sheaf of
meromorphic units whose divisors are supported
on $D$.

\emph{The exact bounded-character subsheaf.}
The period pairing remains (B.4.24.1), pulled
back to $S$. By the base-change lemma
\citep[Lemma 1.2.9, pp.~80--81]{KKN2008Root}
it defines $\mathcal A_S$. More explicitly its
bounded multiplicative-character sheaf is
\[
 \mathcal T_{\log}^{(Y)}
   =\mathcal M_{D_1}^{\mathrm{gp}}
          \times\mathcal M_{D_2}^{\mathrm{gp}}.
 \tag{B.4.27.3}
\]
Indeed the first character is bounded by powers
of $q$, and therefore has zero order on $D_R$.
It may have arbitrary integer order on $D_Q$
and on each $E_i$, where $q$ has positive order.
The second character is bounded by powers of
$q r^{12}$, whose order is positive on every
boundary component, so all its boundary orders
are permitted. These individual conditions
also suffice for every character combination:
if its second coefficient is nonzero, a suitable
signed multiple of the second period bounds
both local normal orders; otherwise powers of
the first period do so on the permitted support.
At a crossing one chooses the bounding powers
simultaneously for its two integer orders.
This proves (B.4.27.3), including at crossings
and at intersections with $D_R$; replacing the
bounded sheaf by the full log-character sheaf
in both coordinates would be incorrect.

\emph{Cohomology of the blown-up bidisc.}
We record the elementary ingredients to fix
the analytic scope. For a point blowup $p$ of
a smooth surface,
$p_*\mathcal O=\mathcal O$ and
$R^j p_*\mathcal O=0$ for $j>0$.
The first identity follows from Hartogs.
For the second, the two standard blowup charts
give a local Laurent calculation: on their
intersection a holomorphic function expands
as a convergent sum of $u^k z^\ell$ with
$k\ge0$. Terms with $\ell\ge0$ extend in
the first chart. Terms with $\ell<0$ extend
in the second, because in coordinates
$z'=1/z$, $u'=uz$ they are
$(u')^k(z')^{k-\ell}$. Cauchy estimates on
smaller coordinate polydiscs preserve convergence
of the two sums. The charts and their
intersection are Stein, so this kills the
first cohomology and there is no higher
\v Cech cohomology. Iterating and using the
Leray sequence over the Stein bidisc gives
\[
 H^j(S,\mathcal O_S)=0\qquad(j>0).
 \tag{B.4.27.4}
\]

The exceptional tree has one $H^2(\ ,\Z)$
generator per rational component and no positive
cohomology in other degrees. Proper base change
and Leray for $\beta$ therefore give
\[
 H^2(S,\Z)\xrightarrow{\ \sim\ }
       \bigoplus_i H^2(E_i,\Z)\simeq\Z^N,
 \qquad H^j(S,\Z)=0\quad(j=1,3,4).
 \tag{B.4.27.5}
\]
Equivalently these groups follow from the usual
tubular plumbing of the exceptional tree.
The exceptional intersection matrix is
unimodular: at each blowup the total-transform
basis adds a vector of square $-1$ orthogonal
to the preceding basis. Passing to the strict
transform basis is an integral unimodular
change of basis. Thus the Chern classes of
$\mathcal O(E_i)$ form a basis of
$H^2(S,\Z)$. The exponential sequence and
(B.4.27.4)--(B.4.27.5) give
\[
 \operatorname{Pic}(S)\simeq\Z^N
       =\langle\mathcal O(E_i)\rangle,
 \qquad
 H^2(S,\mathcal O_S^*)=H^3(S,\mathcal O_S^*)=0.
 \tag{B.4.27.6}
\]

For either $D=D_1$ or $D_2$ the divisor-order
sequence is exact:
\[
 0\longrightarrow\mathcal O_S^*
 \longrightarrow\mathcal M_D^{\mathrm{gp}}
 \longrightarrow
   \bigoplus_{C\subset D}i_{C*}\Z
 \longrightarrow0.
 \tag{B.4.27.7}
\]
The map from global integer orders to the Picard
group sends a component to its divisor line
bundle. It is surjective because both divisors
contain every $E_i$. Every component has
$H^1(C,\Z)=0$: the strict transforms of the
axes are discs and the exceptional curves are
rational. Their only nonzero $H^2$ contributions
are the $N$ exceptional curves. Consequently
the long exact sequence of (B.4.27.7) gives
\[
 H^1(S,\mathcal M_{D_j}^{\mathrm{gp}})=0,
 \qquad
 H^2(S,\mathcal M_{D_j}^{\mathrm{gp}})
       \simeq\bigoplus_i H^2(E_i,\Z)
       \simeq\Z^N\quad(j=1,2).
 \tag{B.4.27.8}
\]
The second isomorphism is induced by the actual
divisor-order map, a fact needed for the next
calculation.

\emph{Injectivity of the period map in degree two.}
The quotient definition of $\mathcal A_S$ gives
\[
 0\longrightarrow Y\longrightarrow
 \mathcal T_{\log}^{(Y)}\longrightarrow
 \mathcal A_S\longrightarrow0,
 \qquad Y=\Z u\oplus\Z\gamma.
 \tag{B.4.27.9}
\]
By (B.4.27.3) and (B.4.27.8), its middle
sheaf has zero $H^1$ and $H^2\simeq\Z^{2N}$.
Using restriction to $E_i$ as the coordinates
in (B.4.27.5), the map
\[
 H^2(S,Y)\longrightarrow
 H^2(S,\mathcal T_{\log}^{(Y)})
\]
is block diagonal. Its block at $E_i$ is the
valuation matrix of the two period columns,
\[
 \begin{pmatrix}
 12m a_i&0\\0&12m a_i+12n b_i
 \end{pmatrix}.
 \tag{B.4.27.10}
\]
Both diagonal entries are positive. Thus this
map is injective over $\Z$, not merely over
$\Q$. The long exact sequence of
(B.4.27.9) now gives (B.4.27.1).

Put $U_S=\beta^{-1}(B\setminus\{0\})$.
The actual torsor on $U_S\simeq B\setminus\{0\}$
is nontrivial by B.4.24, including after the
stated finite coordinate roots. An extending
torsor on $S$ would be trivial by (B.4.27.1),
and its restriction would be trivial as well,
a contradiction. For $N>0$, $U_S$ is the
complement of the exceptional union; for
$N=0$ it is the punctured bidisc. Thus the
case $N=0$ is included as well.

No step supplies a log-group presentation of
an arbitrary analytic cap. In particular a
flattening of its map to the bidisc changes
the base, but does not by itself produce
(B.4.27.9) for that unknown space. The result
therefore removes a specific logarithmic
repair by base blowups. B.4.30--B.4.33 supply the
additional extending-section and global-loop arguments
that exclude an arbitrary proper filling; no such
comparison is assumed in this proposition.
The valuation blocks and the unimodular
intersection matrices for the first 124 marked
point-blowup sequences are independently checked
in \path{code/e_push_log_blowup.py}. This finite
audit supplements the all-length argument above.
\end{proof}
